% Standalone problem document, generated from the adjacent README.md.
% Compile with XeLaTeX. Mathematical content is maintained in Markdown.
\documentclass[11pt,a4paper]{article}
\usepackage[fontsize=10.5pt]{scrextend}
\usepackage[margin=22mm,headheight=15pt,headsep=9mm,footskip=11mm]{geometry}
\usepackage{fontspec}
\setmainfont{texgyrepagella}[Extension=.otf,UprightFont=*-regular,BoldFont=*-bold,ItalicFont=*-italic,BoldItalicFont=*-bolditalic]
\setsansfont{texgyreheros}[Extension=.otf,UprightFont=*-regular,BoldFont=*-bold,ItalicFont=*-italic,BoldItalicFont=*-bolditalic]
\setmonofont{texgyrecursor}[Extension=.otf,UprightFont=*-regular,BoldFont=*-bold,ItalicFont=*-italic,BoldItalicFont=*-bolditalic,Scale=MatchLowercase]
\usepackage{amsmath,amssymb}
\usepackage{unicode-math}
\setmathfont{texgyrepagella-math.otf}
\usepackage{xcolor,microtype,enumitem,needspace,fancyhdr}
\usepackage{bookmark}
\definecolor{ink}{HTML}{17344B}
\definecolor{accent}{HTML}{197D80}
\definecolor{muted}{HTML}{596773}
\hypersetup{colorlinks=true,linkcolor=accent,urlcolor=accent,pdftitle={KE-05 - Spectral-gap-independent
constants for randomized block polynomial
interpolation},pdfauthor={Open Problems in Numerical Linear Algebra}}
\urlstyle{same}
\setlength{\parindent}{0pt}
\setlength{\parskip}{5pt plus 1pt minus 1pt}
\linespread{1.04}
\setlength{\emergencystretch}{3em}
\widowpenalty=10000
\clubpenalty=10000
\displaywidowpenalty=10000
\allowdisplaybreaks[1]
\setlist{nosep,leftmargin=1.5em}
\providecommand{\tightlist}{\setlength{\itemsep}{0pt}\setlength{\parskip}{0pt}}
\providecommand{\pandocbounded}[1]{#1}
\pagestyle{fancy}
\fancyhf{}
\fancyhead[L]{\sffamily\footnotesize\color{muted}OPEN PROBLEMS IN NUMERICAL LINEAR ALGEBRA}
\fancyhead[R]{\sffamily\bfseries\footnotesize\color{accent}KE-05}
\fancyfoot[L]{\parbox[b]{0.9\textwidth}{\sffamily\fontsize{8}{10}\selectfont\color{muted}Editorial ratings; a bounded search cannot certify that a problem remains unsolved.\\Literature check: 2026-09-10}}
\fancyfoot[R]{\sffamily\footnotesize\color{muted}\thepage}
\renewcommand{\headrulewidth}{0.3pt}
\renewcommand{\headrule}{\hbox to\headwidth{\color{accent}\leaders\hrule height \headrulewidth\hfill}}
\makeatletter
\renewcommand{\section}{\Needspace{5\baselineskip}\@startsection{section}{1}{0pt}{12pt plus 2pt minus 2pt}{4pt}{\sffamily\bfseries\large\color{ink}}}
\renewcommand{\subsection}{\Needspace{4\baselineskip}\@startsection{subsection}{2}{0pt}{9pt plus 1pt minus 1pt}{3pt}{\sffamily\bfseries\normalsize\color{ink}}}
\makeatother
\setcounter{secnumdepth}{0}
\begin{document}
{\sffamily\small\color{muted}Randomized and low-rank approximation\par}
\vspace{3pt}
{\raggedright\hyphenpenalty=10000\exhyphenpenalty=10000\sffamily\bfseries\fontsize{21}{24}\selectfont\color{ink}Spectral-gap-independent
constants for randomized block polynomial interpolation\par}
\vspace{7pt}
{\sffamily\small\textbf{Difficulty:} challenging\quad\textbf{Importance:} interesting
to the community\par}
{\sffamily\small\bfseries\color{accent}Status: Open\par}
\vspace{4pt}
{\color{accent}\hrule height 0.8pt}
\vspace{6pt}
\textbf{Rating rationale:} Challenging because almost-sure invertibility
does not yield spectrum-uniform probabilistic constants; community
impact is explaining cluster robustness of block Lanczos.

Fix integers \(b\ge1\), \(d\ge2\). Let
\(\Lambda_1,\ldots,\Lambda_d\in\mathbb R^{b\times b}\) be diagonal with
pairwise disjoint spectra. Draw all entries of
\(\Omega_1,\ldots,\Omega_d\in\mathbb R^{b\times b}\) independently from
\(N(0,1)\) and put \(B_i=\Omega_i^{-1}\Lambda_i\Omega_i\). Let \(a\) and
\(c\) be the smallest and largest diagonal entries in the entire family.

For each \(k\in\{1,\ldots,d\}\), reorder the triples
\((B_i,\Lambda_i,\Omega_i)\) into the order
\((k,1,\ldots,k-1,k+1,\ldots,d)\) and use superscript \((k)\) for that
ordering. In descending order \(i=d,d-1,\ldots,1\), compute \[
S_{i,i}^{(k)}=I_b,\qquad
S_{i,j}^{(k)}=B_i^{(k)}S_{i,j-1}^{(k)}-S_{i,j-1}^{(k)}\widehat B_j^{(k)}
\quad(j=i+1,\ldots,d),
\] \[
\widehat\Omega_i^{(k)}=\Omega_i^{(k)}S_{i,d}^{(k)},\qquad
\widehat B_i^{(k)}=(\widehat\Omega_i^{(k)})^{-1}\Lambda_i^{(k)}\widehat\Omega_i^{(k)}.
\] All norms below are spectral norms. Define \[
\chi_{\rm mono}^{(k)}=
\max_{2\le i\le d}\left\{
1,\frac{\|aI-\widehat B_i^{(k)}\|_2}{\|aI-\Lambda_i^{(k)}\|_2},
\frac{\|cI-\widehat B_i^{(k)}\|_2}{\|cI-\Lambda_i^{(k)}\|_2}\right\},
\] \[
\chi_{\rm coef}^{(k)}=
\|(S_{1,d}^{(k)})^{-1}\|_2^{1/(d-1)}
\min_{\substack{2\le i\le d\ ,\ \lambda\in\sigma(\Lambda_1^{(k)})\\
\eta\in\sigma(\Lambda_i^{(k)})}}|\lambda-\eta|,
\qquad
\chi_{\rm mono}=\max_k\chi_{\rm mono}^{(k)},\quad
\chi_{\rm coef}=\max_k\chi_{\rm coef}^{(k)}.
\] The inverses exist almost surely, as shown in the source. If an
endpoint ratio is \(0/0\), set it to 1; its matrix is necessarily a
scalar matrix and the surrounding maximum already includes 1.

Is the family of random variables \(\chi_{\rm mono}\chi_{\rm coef}\)
uniformly bounded in probability over all admissible diagonal data?
Precisely, for each \(0<\delta<1\), does there exist a finite
\(C(b,d,\delta)\) such that, for every fixed admissible
\((\Lambda_1,\ldots,\Lambda_d)\), \[
\Pr\{\chi_{\rm mono}\chi_{\rm coef}\le C(b,d,\delta)\}\ge1-\delta?
\] The constant must be independent of the eigenvalues and all gaps
within and between their blocks. Probability is taken separately for
each fixed input, not for one draw required to work simultaneously for
every spectrum. This quantifies the source's ``with high probability''
statement with its explicitly allowed failure-probability dependence.

These constants control matrix-polynomial interpolation in the
convergence analysis of randomized small-block Lanczos. The bound would
explain robustness to clusters much larger than the starting block. The
scalar-block case is already bounded deterministically; noncommutativity
is the unresolved obstacle.

\newpage
\subsection{References}\label{references}

N. Shao, \emph{A structural bound for cluster robustness of randomized
small-block Lanczos}, arXiv:2507.10144v2, 30 May 2026, Theorem 1,
equation (12), and Conjecture 1 in §3.2
(\href{https://arxiv.org/html/2507.10144v2}{primary full text}). T. Chen
et al., \emph{Does block size matter in randomized block Krylov low-rank
approximation?}, arXiv:2508.06486, §§4,6
(\href{https://arxiv.org/abs/2508.06486}{primary paper}), is related:
its conjecture changes spectral-gap dependence in output bounds, whereas
this question bounds the intermediate random interpolation constants
themselves.

\subsection{Status check --- 2026-09-08}\label{status-check-2026-09-08}

Latest arXiv abstract/history and May 2026 v2 inspected. The revision
explicitly retains Conjecture 1, including independence of gaps within
and between blocks. Searches for author/title with ``proof'',
``conjecture'', ``cluster robustness'', and 2026 found no resolution.
Generic nonsingularity of the recurrence does not give the claimed
uniform probabilistic bound.

\subsection{Audit --- 2026-09-10}\label{audit-2026-09-10}

Rechecked \href{https://arxiv.org/html/2507.10144v2}{Shao's May 2026
Conjecture 1} and current version history. It still requires
independence of gaps inside and between blocks. Cluster-robustness and
later-author searches found no resolution. The scalar-block baseline
does not resolve noncommuting block interpolation.
\end{document}
